\documentclass[11pt]{article}
\newcommand{\ListaTitulo}{Lista 03: Dispersão: AT, DP, variância e CV}
% Preâmbulo compartilhado: MATF14, Listas de Exercícios 2026
% Usado por ListaNN.tex e GabaritoNN.tex via % Preâmbulo compartilhado: MATF14, Listas de Exercícios 2026
% Usado por ListaNN.tex e GabaritoNN.tex via % Preâmbulo compartilhado: MATF14, Listas de Exercícios 2026
% Usado por ListaNN.tex e GabaritoNN.tex via \input{preamble}

\usepackage[utf8]{inputenc}
\usepackage[T1]{fontenc}
\usepackage[brazil]{babel}
\usepackage{fullpage}
\usepackage{amsmath,amssymb}
\usepackage{enumitem}
\usepackage{xcolor}
\usepackage{fancyhdr}
\usepackage{titlesec}
\usepackage{listings}
\usepackage{hyperref}
\usepackage{booktabs}
\usepackage{graphicx}
\usepackage{tikz}

\definecolor{matf14azul}{HTML}{0D3B54}
\definecolor{matf14dourado}{HTML}{C9971E}

\titleformat{\section}{\normalfont\Large\bfseries\color{matf14azul}}{\thesection}{1em}{}
\titleformat{\subsection}{\normalfont\large\bfseries\color{matf14azul}}{\thesubsection}{1em}{}

\providecommand{\ListaTitulo}{Lista de Exercícios}

\setlength{\headheight}{14pt}
\pagestyle{fancy}
\fancyhf{}
\lhead{\small MATF14: Estatística Econômica I}
\rhead{\small \ListaTitulo}
\cfoot{\thepage}
\renewcommand{\headrulewidth}{0.4pt}

\lstset{
  basicstyle=\ttfamily\small,
  breaklines=true,
  frame=single,
  columns=fullflexible,
  backgroundcolor=\color{gray!8},
  commentstyle=\color{matf14dourado},
  keywordstyle=\color{matf14azul}\bfseries,
  language=R
}

\newcounter{questaocounter}
\newenvironment{questao}[1]{%
  \stepcounter{questaocounter}%
  \bigskip\noindent\textbf{Questão #1.}\ %
}{\par}

\newcommand{\cabecalho}[2]{%
  \begin{center}
    {\Large\bfseries\color{matf14azul} #1}\\[0.3em]
    {\large #2}
  \end{center}
  \vspace{0.5em}
  \noindent\rule{\linewidth}{0.6pt}
  \vspace{0.5em}
}

\newenvironment{solucao}{%
  \par\smallskip\noindent\textit{\color{matf14dourado!80!black}Solução.}\ %
}{\hfill$\blacksquare$\par}


\usepackage[utf8]{inputenc}
\usepackage[T1]{fontenc}
\usepackage[brazil]{babel}
\usepackage{fullpage}
\usepackage{amsmath,amssymb}
\usepackage{enumitem}
\usepackage{xcolor}
\usepackage{fancyhdr}
\usepackage{titlesec}
\usepackage{listings}
\usepackage{hyperref}
\usepackage{booktabs}
\usepackage{graphicx}
\usepackage{tikz}

\definecolor{matf14azul}{HTML}{0D3B54}
\definecolor{matf14dourado}{HTML}{C9971E}

\titleformat{\section}{\normalfont\Large\bfseries\color{matf14azul}}{\thesection}{1em}{}
\titleformat{\subsection}{\normalfont\large\bfseries\color{matf14azul}}{\thesubsection}{1em}{}

\providecommand{\ListaTitulo}{Lista de Exercícios}

\setlength{\headheight}{14pt}
\pagestyle{fancy}
\fancyhf{}
\lhead{\small MATF14: Estatística Econômica I}
\rhead{\small \ListaTitulo}
\cfoot{\thepage}
\renewcommand{\headrulewidth}{0.4pt}

\lstset{
  basicstyle=\ttfamily\small,
  breaklines=true,
  frame=single,
  columns=fullflexible,
  backgroundcolor=\color{gray!8},
  commentstyle=\color{matf14dourado},
  keywordstyle=\color{matf14azul}\bfseries,
  language=R
}

\newcounter{questaocounter}
\newenvironment{questao}[1]{%
  \stepcounter{questaocounter}%
  \bigskip\noindent\textbf{Questão #1.}\ %
}{\par}

\newcommand{\cabecalho}[2]{%
  \begin{center}
    {\Large\bfseries\color{matf14azul} #1}\\[0.3em]
    {\large #2}
  \end{center}
  \vspace{0.5em}
  \noindent\rule{\linewidth}{0.6pt}
  \vspace{0.5em}
}

\newenvironment{solucao}{%
  \par\smallskip\noindent\textit{\color{matf14dourado!80!black}Solução.}\ %
}{\hfill$\blacksquare$\par}


\usepackage[utf8]{inputenc}
\usepackage[T1]{fontenc}
\usepackage[brazil]{babel}
\usepackage{fullpage}
\usepackage{amsmath,amssymb}
\usepackage{enumitem}
\usepackage{xcolor}
\usepackage{fancyhdr}
\usepackage{titlesec}
\usepackage{listings}
\usepackage{hyperref}
\usepackage{booktabs}
\usepackage{graphicx}
\usepackage{tikz}

\definecolor{matf14azul}{HTML}{0D3B54}
\definecolor{matf14dourado}{HTML}{C9971E}

\titleformat{\section}{\normalfont\Large\bfseries\color{matf14azul}}{\thesection}{1em}{}
\titleformat{\subsection}{\normalfont\large\bfseries\color{matf14azul}}{\thesubsection}{1em}{}

\providecommand{\ListaTitulo}{Lista de Exercícios}

\setlength{\headheight}{14pt}
\pagestyle{fancy}
\fancyhf{}
\lhead{\small MATF14: Estatística Econômica I}
\rhead{\small \ListaTitulo}
\cfoot{\thepage}
\renewcommand{\headrulewidth}{0.4pt}

\lstset{
  basicstyle=\ttfamily\small,
  breaklines=true,
  frame=single,
  columns=fullflexible,
  backgroundcolor=\color{gray!8},
  commentstyle=\color{matf14dourado},
  keywordstyle=\color{matf14azul}\bfseries,
  language=R
}

\newcounter{questaocounter}
\newenvironment{questao}[1]{%
  \stepcounter{questaocounter}%
  \bigskip\noindent\textbf{Questão #1.}\ %
}{\par}

\newcommand{\cabecalho}[2]{%
  \begin{center}
    {\Large\bfseries\color{matf14azul} #1}\\[0.3em]
    {\large #2}
  \end{center}
  \vspace{0.5em}
  \noindent\rule{\linewidth}{0.6pt}
  \vspace{0.5em}
}

\newenvironment{solucao}{%
  \par\smallskip\noindent\textit{\color{matf14dourado!80!black}Solução.}\ %
}{\hfill$\blacksquare$\par}


\begin{document}

\cabecalho{Lista de Exercícios 03}{Amplitude total, desvio padrão, variância e coeficiente de variação (Cap. 1, Aula 06)}

\noindent \textit{Justifique todas as respostas. As resoluções são discutidas em sala e nos horários de atendimento; não são publicadas neste material.}

\begin{questao}{1} \textbf{(Cálculo direto: retorno de dois fundos de investimento)}

Os retornos mensais (\%) de dois fundos de investimento, nos últimos 5 meses, foram:
$$
\text{Fundo A: } 1{,}0,\ 1{,}2,\ 0{,}9,\ 1{,}1,\ 0{,}8 \qquad
\text{Fundo B: } -2{,}0,\ 4{,}0,\ 0{,}5,\ 3{,}0,\ -0{,}5
$$

\begin{enumerate}[label=(\alph*)]
  \item Calcule a média de retorno de cada fundo.
  \item Calcule a variância e o desvio padrão (divisor $n$) de cada fundo.
  \item Os dois fundos têm médias parecidas. Um investidor avesso a risco prefere qual fundo?
    Justifique com base no desvio padrão.
\end{enumerate}
\end{questao}

\begin{questao}{2} \textbf{(Coeficiente de variação: comparando variáveis em unidades diferentes)}

Uma pesquisa com 50 empresas registrou o número de funcionários (média $=45$, desvio padrão
$=18$) e o faturamento anual em milhões de reais (média $=12$, desvio padrão $=3$).

\begin{enumerate}[label=(\alph*)]
  \item É correto afirmar que "o número de funcionários é mais disperso que o faturamento, porque
    seu desvio padrão (18) é maior que o do faturamento (3)"? Justifique.
  \item Calcule o $CV$ de cada variável e responda: qual delas é \emph{relativamente} mais
    dispersa em torno de sua própria média?
\end{enumerate}
\end{questao}

\begin{questao}{3} \textbf{(Amplitude total vs. desvio padrão: sensibilidade a outliers)}

Dois times de vendedores têm as seguintes vendas mensais (em unidades vendidas):
$$
\text{Time 1: } 20, 22, 21, 23, 19, 20, 21 \qquad
\text{Time 2: } 20, 21, 20, 22, 19, 21, 90
$$

\begin{enumerate}[label=(\alph*)]
  \item Calcule a amplitude total ($AT$) de cada time.
  \item Calcule o desvio padrão de cada time e compare com o resultado de (a). O que essa
    comparação revela sobre a sensibilidade relativa de $AT$ e de $dp$ a um único valor atípico
    (o "90" do Time 2, um provável erro de digitação ou uma venda excepcional isolada)?
  \item Qual medida de dispersão você recomendaria usar para avaliar a consistência normal das
    vendas do Time 2, dado que o valor 90 pode ser um erro a ser investigado, não descartado sem
    checagem? (Relembre a Aula 05.)
\end{enumerate}
\end{questao}

\begin{questao}{4} \textbf{(Dedução: a variância é sempre não negativa)}

Seja $x_1, \ldots, x_n$ um conjunto de dados quaisquer, com média $\bar x$, e
$\mathrm{Var}(X) = \frac{1}{n}\sum_{i=1}^n (x_i-\bar x)^2$.

\begin{enumerate}[label=(\alph*)]
  \item Mostre que $\mathrm{Var}(X) \ge 0$ sempre, qualquer que seja o conjunto de dados.
  \item Mostre que $\mathrm{Var}(X) = 0$ se, e somente se, todos os valores $x_1=x_2=\cdots=x_n$
    são iguais (isto é, não há dispersão nenhuma). (Dica: uma soma de termos não negativos é zero
    se, e só se, todos os termos são zero.)
\end{enumerate}
\end{questao}

\begin{questao}{5} \textbf{(Dedução: fórmula alternativa da variância)}

Mostre que
$$
\mathrm{Var}(X) = \frac{1}{n}\sum_{i=1}^n (x_i-\bar x)^2 = \left(\frac{1}{n}\sum_{i=1}^n x_i^2\right) - \bar{x}^2,
$$
ou seja, "a variância é a média dos quadrados menos o quadrado da média". (Dica: expanda o
quadrado $(x_i-\bar x)^2 = x_i^2 - 2x_i\bar x + \bar x^2$ dentro do somatório e use que
$\sum_i x_i = n\bar x$.) Essa fórmula alternativa é frequentemente mais rápida para cálculo manual, refaça o item (b) da Questão 1 (Fundo A) usando essa fórmula e confira que obtém o mesmo valor.
\end{questao}

\bigskip
\noindent\textit{A questão a seguir não tem resposta única e não é cobrada em prova no mesmo
formato, é para discussão em grupo ou por escrito, antes da próxima aula.}

\begin{questao}{6 (discussão)} \textbf{(Sem resposta única)}

Um gestor de fundos anuncia dois produtos com o mesmo retorno médio histórico, mas se recusa a
divulgar o desvio padrão de nenhum dos dois "para não confundir o cliente". Escreva, em poucas
linhas, um argumento para convencê-lo de que omitir essa informação pode ser mais prejudicial ao
cliente do que divulgá-la, e proponha uma forma simples de comunicar risco a um público
não técnico, sem recorrer só ao termo "desvio padrão".
\end{questao}


\bigskip
\section*{Exercícios complementares}

\begin{enumerate}[label=\textbf{C\arabic*.}, leftmargin=*]
\item Dois alunos tiveram cinco notas em um curso. Aluno A: 2, 10, 6, 6, 6. Aluno B: 2, 6, 10, 4, 8.
\begin{enumerate}[label=(\alph*)]
  \item Qual é a nota média de cada aluno?
  \item Qual é a amplitude total das notas de cada aluno?
  \item Calcule a variância e o desvio padrão (divisor $n$) das notas de cada aluno. Quem teve maior variabilidade?
\end{enumerate}
\item Temperaturas (°C) em duas cidades, em cinco horários. Cidade A: 24, 26, 32, 28, 25. Cidade B: 20, 26, 36, 28, 25.
\begin{enumerate}[label=(\alph*)]
  \item Mostre que as duas cidades têm a mesma média e o mesmo $IQR$.
  \item Calcule o desvio padrão de cada cidade. Qual é mais variável?
  \item Calcule o coeficiente de variação da cidade B e compare com o das notas do aluno B da questão anterior. O que o $CV$ permite concluir que o desvio padrão não permite?
\end{enumerate}
\item O número de crianças em 20 casas visitadas numa pesquisa foi: 2, 0, 0, 1, 5, 3, 4, 0, 1, 1, 2, 1, 3, 0, 4, 6, 3, 2, 0, 2.
\begin{enumerate}[label=(\alph*)]
  \item Construa a tabela de frequências absolutas.
  \item Usando a tabela, calcule a moda, a mediana e a média.
  \item Usando a tabela, calcule a variância e o desvio padrão.
  \item Use os salários da Companhia MB (Lista 01, questão C3) e calcule a amplitude total, o $IQR$ e o coeficiente de variação dos salários.
\end{enumerate}

\end{enumerate}

\bigskip
\needspace{12\baselineskip}
\section*{Aprofundamento: contas nacionais e dados reais}
\noindent\textit{Exercícios com dados oficiais, ligados à seção de contas nacionais do Capítulo 1 do livro. Use variância com divisor $n$, como em aula.}

\begin{enumerate}[label=\textbf{A\arabic*.}, leftmargin=*]
\item \textbf{(Multiplicadores de insumo-produto.)} Na matriz de insumo-produto, o multiplicador de produção de um setor indica quantos reais
de produção são gerados em toda a economia por R\$ 1,00 a mais de demanda final por produtos desse setor (soma da coluna da matriz inversa de Leontief).
Multiplicadores de 12 setores (IBGE, Matriz de Insumo-Produto 2015):
\begin{center}
\small
\begin{tabular}{lc|lc}
\toprule
Setor & Mult. & Setor & Mult. \\
\midrule
Agropecuária & 1,72 & Comércio & 1,53 \\
Indústrias extrativas & 1,77 & Transporte e armazenagem & 1,84 \\
Indústrias de transformação & 2,15 & Informação e comunicação & 1,64 \\
Eletricidade, gás, água e esgoto & 1,95 & Atividades financeiras & 1,49 \\
Construção & 1,81 & Atividades imobiliárias & 1,11 \\
 & & Outros serviços & 1,53 \\
 & & Administração pública & 1,38 \\
\bottomrule
\end{tabular}
\end{center}
\begin{enumerate}[label=(\alph*)]
  \item Calcule a amplitude total, a média, a variância, o desvio padrão e o coeficiente de variação dos multiplicadores.
  \item Interprete o multiplicador da indústria de transformação. Por que o das atividades imobiliárias é tão próximo de 1?
  \item Um governo quer maximizar o efeito de R\$ 1 bilhão em compras sobre a produção. Os multiplicadores são suficientes para essa decisão? Discuta.
\end{enumerate}

\item \textbf{(Volatilidade dos componentes da demanda.)} Taxas de variação em relação ao trimestre anterior, com ajuste sazonal, do 1º trimestre de 2023 ao
2º trimestre de 2026 (\%; IBGE, Contas Nacionais Trimestrais). Resumos: consumo das famílias, média 0,60 e desvio padrão 0,73; FBCF,
média 0,49 e desvio padrão 2,37; PIB, média 0,69 e desvio padrão 0,50.
\begin{enumerate}[label=(\alph*)]
  \item Calcule o coeficiente de variação de cada série.
  \item O CV da FBCF é quase 5. Por que o CV é uma medida pouco confiável para essas séries? Que medida é mais adequada para comparar a volatilidade
    delas? (Dica: todas estão na mesma unidade, pontos percentuais.)
  \item Por que o investimento costuma ser muito mais volátil que o consumo? Relacione com a sua resposta a (b).
\end{enumerate}

\end{enumerate}

\end{document}
